\chapter{OPERATOR PADA RUANG HILBERT}

\section{Pemetaan Linear Invertibel dan Isometri}

\begin{defn} Suatu pemetaan linear $T\colon H \sto K$ antara dua ruang hasil kali dalam dikatakan
 \index{hasil kali dalam!mempertahankan}%
\df{mempertahankan hasil kali dalam} jika $\langle Tx, Ty \rangle = \langle x, y \rangle$ untuk semua $x$, $y \in H$.
\end{defn}

\begin{prop}\label{prop_isometry_equiv} Misalkan $T\colon H \sto K$ suatu pemetaan linear antara ruang-ruang hasil kali dalam. Maka
pernyataan-pernyataan berikut ekuivalen:
 \begin{enumerate}[label=\rm{(\alph*)}]
   \item $T$ mempertahankan hasil kali dalam;
   \item $T$ mempertahankan norma; dan
   \item $T$ merupakan isometri.
 \end{enumerate}
\end{prop}

Sama seperti pada ruang linear bernorma (lihat bagian paling akhir dari seksi~\ref{sec_bdd_lin_maps}),
terdapat (sekurang-kurangnya) dua kategori penting yang objek-objeknya adalah ruang Hilbert:
 \begin{enumerate}
  \item[(a)] $\cat{HSp} =$ ruang Hilbert dan pemetaan linear terbatas, dan
   \index{HSp@$\cat{HSp}$!kategori}%
   \index{kategori!$\cat{HSp}$ sebagai suatu}%
  \item[(b)] $\cat{HSp_1} =$ ruang Hilbert dan kontraksi linear.
   \index{HSp1@$\cat{HSp_1}$!kategori}%
   \index{kategori!$\cat{HSp_1}$ sebagai suatu}%
 \end{enumerate}
Dalam kategori apa pun, lazim untuk mendefinisikan \emph{isomorfisme} sebagai \emph{morfisme invertibel}. Karena hampir selalu
 \index{kategori!topologis}%
kategori (topologis) ruang Hilbert dan pemetaan linear terbataslah yang kita jumpai
dalam catatan ini, tampaknya masuk akal untuk menerapkan kata ``isomorfisme'' pada
isomorfisme dalam kategori ini---dengan kata lain, pada pemetaan invertibel; sedangkan isomorfisme
dalam kategori yang lebih ketat, yaitu
 \index{kategori!geometris}%
kategori (geometris) ruang Hilbert dan kontraksi linear, dapat secara wajar
disebut \emph{isomorfisme isometrik}. Namun, ini \emph{bukan} konvensi yang dipakai. Ketika kebanyakan matematikawan memikirkan
``isomorfisme ruang Hilbert'', yang mereka maksud adalah pemetaan yang mempertahankan seluruh struktur ruang Hilbert---termasuk
hasil kali dalam. Dengan demikian (menggunakan~\ref{prop_isometry_equiv}), kata
 \index{isomorfisme!antara ruang Hilbert}%
 \index{isomorfisme!dalam kategori $\cat{HSp}$}%
 \index{HSp@$\cat{HSp}$!isomorfisme dalam}%
 \index{isomorfisme!dalam kategori $\cat{HSp_1}$}%
 \index{HSp1@$\cat{HSp_1}$!isomorfisme dalam}%
``isomorfisme'', apabila diterapkan pada pemetaan antara ruang-ruang Hilbert, di hampir semua tempat berarti \emph{isomorfisme isometrik}.
Akibatnya, isomorfisme dalam kategori yang lebih umum $\cat{HSp}$ disebut \emph{pemetaan (linear terbatas) invertibel}.
Ingat pula: Kita mencadangkan kata
 \index{operator}%
 \index{konvensi!semua operator terbatas dan linear}%
``operator'' bagi pemetaan linear terbatas dari suatu ruang \emph{ke ruang itu sendiri}.

\begin{exam} Definisikan suatu operator $T$ pada ruang Hilbert $H = L_2([0,\infty))$ dengan
   \[ Tf(t) = \begin{cases}   f(t-1),  &\text{ jika $t \ge 1$} \\
                                   0,  &\text{ jika $0 \le t < 1$.}
           \end{cases} \]
Maka $T$ merupakan isometri, tetapi bukan isomorfisme isometrik. Sebaliknya, jika
$H = L_2(\R)$ dan $T$ memetakan setiap $f \in H$ ke fungsi $t \mapsto f(t-1)$, maka $T$ merupakan isomorfisme isometrik.
\end{exam}

\begin{exam} Misalkan $H$ ruang Hilbert yang didefinisikan dalam contoh~\ref{exam_Hsp_abs_cont}. Maka pemetaan diferensiasi
$D\colon f \mapsto f\,'$ dari $H$ ke $L_2([0,1])$ merupakan isomorfisme isometrik. Apa inversnya?
\end{exam}

\begin{exam} Misalkan $\{(X_i,\mu_i) \colon i \in I\}$ suatu keluarga ruang ukur $\sigma$-hingga. Jadikan gabungan saling lepas
$\D X = \biguplus_{i \in I}X_i$ sebagai ruang ukur $(X,\mu)$ dengan cara biasa. Maka $L_2(X,\mu)$ isomorfik secara isometrik
dengan jumlah langsung ruang-ruang Hilbert $L_2(X_i,\mu_i)$.
\end{exam}

\begin{exam} Misalkan $\mu$ ukuran Lebesgue pada $[-\frac\pi2, \frac\pi2]$ dan $\nu$ ukuran Lebesgue pada $\R$.
Misalkan $(Uf)(x) = \dfrac{f(\arctan x)}{\sqrt{1+x^2}}$ untuk semua $f \in L_2(\mu)$ dan semua $x \in \R$. Maka $U$ merupakan
isomorfisme isometrik antara $L_2([-\frac\pi2, \frac\pi2], \mu)$ dan $L_2(\R, \nu)$.
\end{exam}

\begin{prop} Setiap surjeksi yang mempertahankan hasil kali dalam dari suatu ruang Hilbert ke ruang Hilbert lainnya secara otomatis linear.
\end{prop}

\begin{prop}\label{prop_norm_lin_map} Jika $T\colon H \sto K$ merupakan pemetaan linear terbatas antara ruang-ruang Hilbert tak nol, maka
   \[ \norm T = \sup\{\abs{\langle Tx,y \rangle}\colon \norm x = \norm y = 1\}\,. \]
\end{prop}

\begin{proof}[\emph{Petunjuk pembuktian}] Tunjukkan terlebih dahulu bahwa $\norm x = \sup\{\abs{\langle x,y \rangle}\colon \norm y = 1\}$. \ns
\end{proof}

\begin{defn} Suatu
 \index{kurva}%
\df{kurva} dalam ruang Hilbert $H$ adalah pemetaan kontinu dari $[0,1]$ ke~$H$. Suatu kurva dikatakan
 \index{sederhana!kurva}%
\df{sederhana} jika kurva itu injektif. Jika $c$ suatu kurva dan $0 \le a < b \le 1$, maka
 \index{tali busur}%
\df{tali busur} dari $c$ yang bersesuaian dengan interval $[a,b]$ adalah vektor $c(b) - c(a)$. Dua tali busur dikatakan
\df{tak bertumpang tindih} jika interval parameter yang terkait dengannya paling banyak memiliki satu titik ujung bersama.
\end{defn}

Dalam \cite{Halmos:1982}, Halmos menggambarkan betapa luasnya ruang Hilbert berdimensi tak hingga melalui
contoh elegan berikut.

\begin{exam} Dalam setiap ruang Hilbert berdimensi tak hingga terdapat kurva sederhana yang berbelok membentuk sudut siku-siku
di setiap titik, dalam arti bahwa setiap pasangan tali busur yang tak bertumpang tindih saling tegak lurus.
\end{exam}

\begin{proof}[\emph{Petunjuk pembuktian}] Tinjau fungsi-fungsi karakteristik dalam $L_2([0,1])$. \ns
\end{proof}
























\section{Operator dan Adjoinnya}

\begin{prop}\label{C0635103} Misalkan $S$, $T \in \ofml B(H,K)$, dengan $H$ dan $K$ ruang
Hilbert. Jika $\langle Sx,y \rangle = \langle Tx,y \rangle$ untuk setiap $x \in H$ dan $y \in
K$, maka $S = T$.
\end{prop}

\begin{defn} Misalkan $T$ suatu operator pada ruang Hilbert~$H$. Maka $Q_T$, yaitu
 \index{bentuk kuadratik}%
 \index{bentuk!kuadratik}%
 \index{Q@$Q_T$ (bentuk kuadratik yang terkait dengan~$T$)}%
\df{bentuk kuadratik yang terkait dengan~$T$}, adalah fungsi bernilai skalar yang didefinisikan oleh
   \[ Q_T(x) := \langle Tx,x \rangle \]
untuk setiap $x \in H$.
\end{defn}

Suatu operator pada ruang Hilbert kompleks adalah operator nol jika dan hanya jika bentuk kuadratik yang terkait dengannya adalah fungsi nol.

\begin{prop}\label{C0635105} Jika $H$ adalah ruang Hilbert \emph{kompleks} dan $T \in \ofml B(H)$, maka
$T = \vc 0$ jika dan hanya jika $Q_T = 0$.
\end{prop}

\begin{proof}[\emph{Petunjuk pembuktian}] Dalam hipotesis $Q_T(z) = 0$ untuk setiap $z \in H$, gantilah $z$ mula-mula dengan $y + x$
dan kemudian dengan $y + ix$.  \ns
\end{proof}

Proposisi sebelumnya merupakan salah satu dari sedikit proposisi yang akan kita jumpai yang tidak berlaku sekaligus untuk ruang Hilbert real dan kompleks.

\begin{exam} Proposisi~\ref{C0635105} tidak benar jika dalam hipotesis kata ``kompleks'' diganti dengan ``real''.
\end{exam}

\begin{defn} Suatu fungsi $T\colon V \sto W$ antara ruang vektor kompleks disebut
 \index{konjugat!linear}%
 \index{linear!konjugat}%
\df{linear konjugat} jika fungsi itu aditif ($T(x+y) = Tx + Ty$) dan memenuhi $T(\alpha x) =
\conj\alpha\, Tx$ untuk semua $x \in V$ dan $\alpha \in \C$. Suatu fungsi bernilai kompleks dengan dua
variabel $\phi\colon V \times W \sto \C$ disebut
 \index{fungsional seskuilinear}%
 \index{fungsional!seskuilinear}%
\df{fungsional seskuilinear} jika fungsi itu linear dalam variabel pertamanya dan linear konjugat
dalam variabel keduanya.

Jika $H$ dan $K$ adalah ruang linear bernorma, suatu fungsional seskuilinear $\phi$ pada $H \times K$
disebut
 \index{terbatas!fungsional seskuilinear}%
 \index{fungsional seskuilinear!terbatas}%
 \index{fungsional!seskuilinear terbatas}%
\df{terbatas} jika terdapat suatu konstanta $M > 0$ sedemikian sehingga
    \[ \abs{\phi(x,y)} \le M\norm x\norm y \]
untuk semua $x \in H$ dan $y \in K$.
\end{defn}

\begin{prop} Jika $\phi \colon H \oplus K \sto \C$ adalah fungsional seskuilinear terbatas pada
jumlah langsung dua ruang Hilbert tak nol, maka bilangan-bilangan berikut semuanya ada dan sama:
 \begin{enumerate}[label=\rm{(\alph*)}]
    \item $\sup\{\abs{\phi(x,y)}\colon \norm x \le 1, \norm y \le 1\}$
 \vskip 3 pt
    \item $\sup\{\abs{\phi(x,y)}\colon \norm x = \norm y = 1\}$
 \vskip 3 pt
    \item $\sup\left\{\dfrac{\abs{\phi(x,y)}}{\norm x\,\norm y}\colon x,y \ne 0\right\}$
 \vskip 3 pt
    \item $\inf\{M > 0\colon \abs{\phi(x,y)} \le M\norm x\,\norm y
                    \text{ untuk semua $x \in H$, $y \in K$}\}$.
 \end{enumerate}
\end{prop}

\begin{proof}[\emph{Petunjuk pembuktian}] Pembuktiannya hampir sama persis dengan hasil yang bersesuaian
untuk pemetaan linear (lihat~\ref{C023221}). \ns
\end{proof}

\begin{defn} Misalkan $\phi \colon H \oplus K \sto \C$ fungsional seskuilinear terbatas pada
jumlah langsung dua ruang Hilbert tak nol. Kita mendefinisikan $\norm \phi$, yaitu
 \index{norma!fungsional seskuilinear terbatas}%
\emph{norma} dari $\phi$, sebagai salah satu di antara ekspresi-ekspresi (yang sama) dalam hasil sebelumnya.
\end{defn}

\begin{prop} Misalkan $T\colon H \sto K$ pemetaan linear terbatas antara ruang Hilbert tak nol. Maka
   \[ \phi \colon H \oplus K \sto \C\colon (x,y) \mapsto \langle Tx,y \rangle \]
adalah fungsional seskuilinear terbatas pada $H \oplus K$ dan $\norm\phi = \norm T$.
\end{prop}

\begin{prop} Misalkan $\phi \colon H \times K \sto \C$ fungsional seskuilinear terbatas pada
hasil kali dua ruang Hilbert tak nol. Maka terdapat pemetaan linear terbatas yang unik $T \in \ofml B(H,K)$ dan
$S \in \ofml B(K,H)$ sedemikian sehingga
    \[ \phi(x,y) = \langle Tx,y \rangle = \langle x,Sy \rangle \]
untuk semua $x \in H$ dan $y \in K$. Selain itu, $\norm T = \norm S = \norm\phi$.
\end{prop}

\begin{proof}[\emph{Petunjuk pembuktian}] Tunjukkan bahwa untuk setiap $x \in H$, pemetaan
$y \mapsto \conj{\phi(x,y)}$ adalah fungsional linear terbatas pada~$K$. \ns
\end{proof}

\begin{defn}\label{def000926} Misalkan $T\colon H \sto K$ pemetaan linear terbatas antara ruang-ruang Hilbert. Pemetaan
$(x,y) \mapsto \langle Tx,y \rangle$ dari $H \oplus K$ ke dalam $\C$ adalah fungsional
seskuilinear terbatas. Berdasarkan proposisi sebelumnya, terdapat pemetaan linear terbatas yang
unik
 \index{adjoin!pemetaan linear!antara ruang Hilbert}%
 \index{linear!pemetaan!adjoin}%
 \index{<unopurstar@$T^*$ (adjoin ruang Hilbert)}%
$T^*\colon K \sto H$ yang disebut \df{adjoin} dari $T$ sedemikian sehingga
   \[ \langle Tx,y \rangle = \langle x,T^*y \rangle \]
untuk semua $x \in H$ dan $y \in K$.
\end{defn}

\begin{exam}\label{C063526} Ingat dari Contoh~\ref{exam150013} bahwa keluarga $l_2$ dari semua
barisan bilangan kompleks yang kuadratnya dapat dijumlahkan merupakan ruang Hilbert. Misalkan
    \[ S\colon l_2 \sto l_2\colon (x_1,x_2,x_3,\dots) \mapsto (0,x_1,x_2,\dots)\,. \]
Maka $S$ adalah operator pada $l_2$, yang disebut
 \index{geser unilateral}%
 \index{S@$S$ (operator geser unilateral)}%
 \index{operator!geser unilateral}%
\df{operator geser unilateral}, dan $\norm S = 1$. Adjoin $S^*$ dari operator geser
unilateral
 \index{geser unilateral!adjoin}%
 \index{adjoin!operator geser unilateral}%
 \index{operator!geser unilateral!adjoin}%
diberikan oleh
  \[ S^*\colon l_2 \sto l_2\colon (x_1,x_2,x_3,\dots) \mapsto (x_2,x_3,x_4,\dots). \]
\end{exam}

\begin{exam} Ambil $\bigl(e^n\bigr)_{n=1}^\infty$ sebagai basis ortonormal bagi ruang Hilbert separabel~$H$ dan
$\bigl(\alpha_n\bigr)_{n=1}^\infty$ barisan terbatas bilangan kompleks. Misalkan $Te^n = \alpha_n e^n$ untuk setiap~$n$.
Maka $T$ dapat diperluas secara unik menjadi operator pada~$H$. Berapakah norma operator ini? Apa adjoinnya?
\end{exam}

\begin{defn} Operator yang dibahas dalam contoh sebelumnya dikenal sebagai
 \index{operator!diagonal}%
 \index{diagonal!operator}%
\df{operator diagonal}. Notasi yang lazim untuk operator ini adalah
 \index{diag@$\diag(\alpha_1,\alpha_2,\dots)$ (operator diagonal)}%
$\diag(\alpha_1,\alpha_2,\dots)$.
\end{defn}

Dalam kategori ruang hasil kali dalam dan pemetaan linear terbatas, baik kernel maupun jangkauan suatu morfisme merupakan objek dalam kategori tersebut.
Namun, seperti ditunjukkan oleh contoh berikut, hal ini tidak harus berlaku untuk ruang Hilbert. Walaupun operator pada ruang Hilbert selalu mempunyai kernel
yang sendiri merupakan ruang Hilbert, jangkauannya mungkin tidak tertutup dan akibatnya tidak lengkap.

\begin{exam}\label{exam_ran_nonclosed} Jangkauan operator diagonal $\diag\bigl(1,\frac12,\frac13, \dots\bigr)$ pada
 \index{tertutup!jangkauan!operator tanpa}%
 \index{operator!tanpa jangkauan tertutup}%
ruang Hilbert $l_2$ tidak tertutup.
\end{exam}

\begin{exam}\label{C063527} Misalkan $(S,\fml A, \mu)$ ruang ukur positif yang sigma-hingga dan $L_2(S,\mu)$
ruang Hilbert dari semua (kelas ekuivalensi) fungsi bernilai kompleks pada $S$
yang terintegralkan kuadrat terhadap~$\mu$. Misalkan $\phi$ fungsi bernilai kompleks yang
terbatas secara esensial dan terukur terhadap $\mu$ pada~$S$. Definisikan
 \index{M@$M_\phi$ (operator perkalian)}%
$M_\phi$ pada $L_2(S,\mu)$ dengan $M_\phi(f) := \phi f$. Maka $M_\phi$ adalah operator pada $L_2(S,\mu)$; operator ini disebut
 \index{operator perkalian!pada $L_2(S,\mu)$}%
 \index{operator!perkalian!pada $L_2(S,\mu)$}%
\df{operator perkalian}. Normanya diberikan oleh $\norm{M_\phi} = \norm\phi_\infty$ dan
 \index{operator perkalian!adjoin}%
 \index{operator!perkalian!adjoin}%
 \index{adjoin!operator perkalian}%
adjoinnya oleh ${M_\phi}^*~=~M_{\conj\phi}$.
\end{exam}

Akan berguna untuk sedikit memperluas Definisi~\ref{0001612}.

\begin{defn} Operator $S$ dan $T$ masing-masing pada ruang Hilbert $H$ dan $K$ disebut
 \index{uniter!ekuivalensi}%
 \index{ekuivalen!secara uniter}%
\df{ekuivalen secara uniter} jika terdapat isomorfisme isometrik $U\colon H \sto K$ sedemikian sehingga $T = USU^{-1}$.

\end{defn}

\begin{exam} Misalkan $\mu$ ukuran Lebesgue pada $[-\frac\pi2, \frac\pi2]$ dan $\nu$ ukuran Lebesgue pada $\R$.
Misalkan $\phi(x) = x$ untuk $x \in [-\frac\pi2, \frac\pi2]$ dan $\psi(x) = \arctan x$ untuk $x \in \R$. Maka operator
perkalian $M_\phi$ dan $M_\psi$ (masing-masing pada $L_2([-\frac\pi2, \frac\pi2], \mu)$ dan $L_2(\R, \nu)$)
ekuivalen secara uniter.
\end{exam}

\begin{exam} Misalkan $M_\phi$ operator perkalian pada $L_2(S,\mu)$, dengan $(S,\fml A,\mu)$
ruang ukuran $\sigma$-hingga seperti dalam Contoh~\ref{C063527}. Maka $M_\phi$ idempoten jika dan hanya jika
$\phi$ adalah fungsi karakteristik dari suatu himpunan dalam~$\fml A$.
\end{exam}

Kita dapat memandang
 \index{spektral!teorema}%
\df{teorema spektral} dalam aljabar linear dasar sebagai pernyataan berikut: \textbf{Setiap operator normal pada ruang hasil kali dalam kompleks
berdimensi hingga ekuivalen secara uniter dengan suatu operator perkalian.} Yang benar-benar
menakjubkan adalah bagaimana pernyataan ini digeneralisasi ke ruang Hilbert berdimensi tak hingga---pokok bahasan
yang akan dibicarakan jauh kemudian.

\begin{exer} Misalkan $\N_3 = \{1,2,3\}$ dan $\mu$ ukuran pencacahan pada $\N_3$.
 \begin{enumerate}
  \item[(a)] Identifikasilah ruang Hilbert~$L_2(\N_3,\mu)$.
  \item[(b)] Identifikasilah operator perkalian pada~$L_2(\N_3,\mu)$.
  \item[(c)] Misalkan $T$ operator linear pada $\C^3$ yang representasi
matriksnya adalah
   \[ \begin{bmatrix}
          2  &  0  &  1 \\
          0  &  2  & -1 \\
          1  & -1  &  1
      \end{bmatrix}.\]
Gunakan operator $T$ untuk mengilustrasikan rumusan teorema spektral sebelumnya.
 \end{enumerate}
\end{exer}

\begin{exam}\label{exam_int_op} Misalkan $L_2 = L_2([0,1],\lambda)$ ruang Hilbert dari semua (kelas ekuivalensi)
fungsi bernilai kompleks pada $[0,1]$ yang terintegralkan kuadrat terhadap ukuran Lebesgue~$\lambda$.
Jika $k\colon [0,1] \times [0,1] \sto \C$ fungsi terukur Borel yang terbatas, definisikan $K$ pada $L_2$ dengan
   \[ (Kf)(x) = \int_0^1 k(x,y)f(y)\,dy\,. \]
Maka $K$ adalah
 \index{integral!operator}%
 \index{operator!integral}%
\df{operator integral} dan
 \index{kernel!operator integral}%
 \index{integral!operator!kernel}%
\df{kernel} dari operator tersebut adalah fungsi~$k$. (Ini merupakan penggunaan lain dari kata ``kernel''; penggunaan ini sama sekali tidak berkaitan
dengan penggunaan kata yang lebih umum: $\ker K = K^{\gets}(\{\vc 0\})$---lihat Definisi~\ref{defn_op_vsp}\,).
Adjoin $K^*$ dari operator $K$ juga merupakan operator integral pada $L_2$ dan
 \index{adjoin!operator integral}%
 \index{integral!operator!adjoin}%
 \index{operator!integral!adjoin}%
kernelnya adalah fungsi $k^*$ yang didefinisikan pada $[0,1] \times [0,1]$ oleh $k^*(x,y) = \conj{k(y,x)}$.
\end{exam}

\begin{proof}[\emph{Petunjuk pembuktian}] Verifikasilah bahwa untuk $f \in L_2$
   \[ \abs{(Kf)(x)} \le \norm{k}_\infty^{1/2}\left[ \int_0^1 \abs{k(x,y)}\abs{f(y)}^2\,dy\right]^{1/2}.\]
\ns
\end{proof}

\begin{exam}\label{000319} Misalkan $H = L_2([0,1])$ ruang Hilbert real dari semua (kelas ekuivalensi) fungsi
bernilai real pada $[0,1]$ yang terintegralkan kuadrat terhadap ukuran Lebesgue. Definisikan $V$ pada $H$ dengan
   \begin{equation}\label{000319i}
       Vf(x) = \int_0^x f(t)\,dt\,.
   \end{equation}
Maka $V$ merupakan operator pada~$H$. Operator ini adalah
 \index{operator Volterra}%
 \index{operator!Volterra}%
\df{operator Volterra} dan merupakan contoh operator integral. (Apa kernel~$k$\, untuk operator ini?)
\end{exam}

\begin{exer} Misalkan $V$ adalah \emph{operator Volterra} yang didefinisikan di atas.
 \begin{enumerate}
  \item[(a)] Tunjukkan bahwa $V$ injektif.
  \item[(b)] Hitunglah $V^*$.
  \item[(c)] Tentukan $V + V^*$ beserta jangkauannya.
 \end{enumerate}
\end{exer}

\begin{prop}\label{prop_adj_invol} Jika $S$ dan $T$ adalah operator pada ruang Hilbert kompleks~$H$ dan $\alpha \in \C$, maka
 \begin{enumerate}[label=\rm{(\alph*)}]
  \item $(S + T)^* = S^* + T^*$;
  \item $(\alpha T)^* = \conj\alpha T^*$;
  \item $T^{**} = T$; dan
  \item $(TS)^* = S^*T^*$.
 \end{enumerate}
\end{prop}

\begin{prop}\label{C063544} Misalkan $T$ operator pada ruang Hilbert~$H$. Maka
   \begin{enumerate}[label=\rm{(\alph*)}]
     \item $\norm{T^*} = \norm T$ dan
     \item $\norm{T^*T} = {\norm T}^2$.
   \end{enumerate}
\end{prop}

\begin{proof}[\emph{Petunjuk pembuktian}] Perhatikan bahwa karena $T = T^{**}$, untuk membuktikan (a) cukup ditunjukkan bahwa $\norm T \le \norm{T^*}$.
Untuk itu, amati bahwa ${\norm{Tx}}^2 = \langle T^*Tx,x \rangle \le \norm{T^*}\norm T$ apabila $x \in H$ dan $\norm x \le 1$. Dari sini
simpulkan bahwa ${\norm T}^2 \le \norm{T^*T} \le \norm{T^*}\norm T$.  \ns
\end{proof}

Syarat (b) yang tampak cukup biasa dalam proposisi sebelumnya kelak ternyata menjadi \emph{sifat} yang sangat penting
ketika aljabar-$C^*$ pertama kali diperkenalkan dalam Bab~\futurexref{7}{chap_cpt_ops}.

\begin{prop}\label{prop_ker_adj_perp_ran_op} Jika $T$ adalah operator pada ruang Hilbert, maka
 \begin{enumerate}[label=\rm{(\alph*)}]
   \item $\ker T^* = (\ran T)^\perp$,
   \item $\clo{\ran T^*} = (\ker T)^\perp$,
   \item $\ker T = (\ran T^*)^\perp$, dan
   \item $\clo{\ran T} = (\ker T^*)^\perp$.
 \end{enumerate}
\end{prop}

Bandingkan hasil sebelumnya dengan Teorema~\ref{00016025}.

\begin{defn}\label{def_bdd_away_zero} Suatu pemetaan linear terbatas $T \colon V \sto W$ antara ruang linear bernorma disebut
 \index{terbatas!jauh dari nol}%
 \index{nol!terbatas jauh dari}%
\df{terbatas jauh dari nol} (atau
 \index{terbatas!dari bawah}%
\df{terbatas dari bawah}) jika terdapat suatu bilangan $\delta > 0$ sedemikian sehingga $\norm{Tx} \ge \delta\norm{x}$ untuk setiap $x \in V$.
\end{defn}

\begin{exam} Jelas bahwa jika suatu pemetaan linear terbatas memiliki sifat terbatas jauh dari nol, maka pemetaan itu injektif. Operator
$T\colon x \mapsto (x_1, \frac12 x_2, \frac13 x_3, \dots)$ yang didefinisikan pada ruang Hilbert $l_2$ menunjukkan bahwa sifat terbatas
jauh dari nol merupakan syarat yang sungguh lebih kuat daripada sifat injektif.
\end{exam}

\begin{prop}\label{prop_bafz_closed_rng} Setiap pemetaan linear terbatas antara ruang Hilbert yang terbatas jauh dari
nol memiliki jangkauan tertutup.
\end{prop}

\begin{prop}\label{prop_nasc_op_invert} Suatu operator pada ruang Hilbert invertibel jika dan hanya jika
operator itu terbatas jauh dari nol dan memiliki jangkauan padat.
\end{prop}

\begin{prop} Pasangan pemetaan $H \mapsto H$ dan $T \mapsto T^*$ yang memetakan setiap ruang Hilbert ke
 \index{adjoin!sebagai funktor pada $\cat{HSp}$}%
 \index{funktor!adjoin}%
dirinya sendiri dan setiap pemetaan linear terbatas antara ruang Hilbert ke adjoinnya merupakan funktor kontravarian
dari kategori $\cat{HSp}$ yang terdiri atas ruang Hilbert dan pemetaan linear terbatas ke kategori itu sendiri.
 \index{HSp@$\cat{HSp}$!funktor adjoin pada}%
\end{prop}

\begin{prop} Misalkan $H$ ruang Hilbert, $T \in \ofml B(H)$, $M = H \oplus \{0\}$, dan $N$ adalah graf dari~$T$.
 \begin{enumerate}[label=\rm{(\alph*)}]
  \item Maka $N$ adalah subruang linear tertutup dari $H \oplus H$.
  \item $T$ injektif jika dan hanya jika $M \cap N = \{(0,0)\}$.
  \item $\ran T$ padat dalam $H$ jika dan hanya jika $M + N$ padat dalam $H \oplus H$.
  \item $T$ surjektif jika dan hanya jika $M + N = H \oplus H$.
 \end{enumerate}
\end{prop}

\begin{exam} Terdapat ruang Hilbert $H$ dengan subruang $M$ dan $N$ sedemikian sehingga $M + N$ tidak tertutup.
\end{exam}

\begin{prop} Misalkan $H$ dan $K$ ruang Hilbert dan $V$ adalah subruang vektor dari~$H$. Setiap pemetaan linear terbatas
$T\colon V \sto K$ dapat diperluas menjadi pemetaan linear terbatas dari $\clo V$, yaitu tutupan $V$, tanpa memperbesar normanya.
\end{prop}

\begin{exer} Misalkan $\fml E$ adalah basis ortonormal bagi ruang Hilbert~$H$. Bahas transformasi linear $T$ pada $H$
sedemikian sehingga $T(e) = 0$ untuk setiap $e \in \fml E$. Berikan contoh tak trivial.
\end{exer}






















\section{Aljabar dengan Involusi}
\begin{defn}\label{00109}  Suatu
 \index{involusi}%
\df{involusi} pada suatu aljabar $A$ adalah suatu pemetaan $x \mapsto x^\ast$ dari $A$ ke $A$ yang memenuhi
  \begin{enumerate}
    \item[(a)] $(x + y)^\ast = x^\ast + y^\ast$,
    \item[(b)] $(\alpha x)^* = \conj\alpha\, x^*$,
    \item[(c)] $x^{\ast\ast} = x$, dan
    \item[(d)] $(xy)^\ast = y^\ast x^\ast$
  \end{enumerate}
untuk semua $x$, $y \in A$, dan $\alpha \in \C$.  Suatu aljabar yang padanya telah didefinisikan suatu involusi adalah
 \index{star@$*\,$-aljabar}%
 \index{<star@$*\,$-aljabar}%
 \index{aljabar!dengan involusi}%
 \index{aljabar!$*\,$-}%
\df{aljabar-$*\,$} (dibaca ``aljabar bintang''). Jika $a$ adalah elemen dari suatu aljabar-$*\,$, maka $a^*$ disebut
 \index{adjoin}%
\df{adjoin} dari~$a$.
\end{defn}

\begin{exam} Dalam aljabar $\C$ dari bilangan kompleks, pemetaan $z \mapsto \conj z$ yang memetakan suatu
bilangan ke konjugat kompleksnya merupakan suatu involusi.
\end{exam}

\begin{exam}\label{0010902} Pemetaan yang memetakan suatu matriks $n \times n$ ke transpos konjugatnya
merupakan suatu involusi pada aljabar beridentitas~$M_n$.
\end{exam}

\begin{exam} Misalkan $X$ suatu ruang Hausdorff kompak.  Pemetaan $f \mapsto \conj f$ yang memetakan suatu
fungsi ke konjugat kompleksnya merupakan suatu involusi dalam aljabar~$\fml C(X)$.
\end{exam}

\begin{exam} Pemetaan $T \mapsto T^*$ yang memetakan suatu operator ruang Hilbert ke adjoinnya merupakan suatu
involusi dalam aljabar $\ofml B(H)$ (lihat proposisi~\ref{prop_adj_invol}).
\end{exam}

\begin{prop} Misalkan $a$ dan $b$ elemen dari suatu aljabar-$*\,$.  Maka $a$ berkomutasi dengan $b$
jika dan hanya jika $a^*$ berkomutasi dengan $b^*$.
\end{prop}

\begin{prop} Dalam suatu aljabar-$*\,$ beridentitas, $\vc 1^* = \vc 1$.
\end{prop}

\begin{prop}\label{00109125} Jika suatu aljabar-$*\,$ $A$ memiliki identitas perkalian kiri $e$,
maka $A$ beridentitas dan $e = \vc 1_A$.
\end{prop}

\begin{prop}\label{00109128fa} Misalkan $a$ elemen dari suatu aljabar-$*\,$ beridentitas.  Maka $a^*$ invertibel jika dan hanya jika $a$ invertibel.
 \index{invers!dari suatu adjoin}%
 \index{invers!adjoin dari suatu}%
 \index{adjoin!invers dari suatu}%
 \index{adjoin!dari suatu invers}%
Dan ketika $a$ invertibel, berlaku
   \[ (a^*)^{-1} = \bigl(a^{-1}\bigr)^*\,. \]
\end{prop}

\begin{defn} Suatu homomorfisme aljabar $\phi$ di antara aljabar-$*\,$ yang mempertahankan
involusi (yaitu, sedemikian sehingga $\phi(a^*) = (\phi(a))^*$) adalah suatu
 \index{<star@homomorfisme-$*\,$}%
 \index{star@homomorfisme-$*\,$}%
 \index{homomorfisme!$*\,$-}%
\df{homomorfisme-$*\,$} (dibaca ``homomorfisme bintang''). Suatu homomorfisme-$*\,$ $\phi\colon A \sto B$ di antara aljabar beridentitas disebut
 \index{beridentitas!homomorfisme-$*\,$}%
 \index{morfisme bijektif!dapat dibalik!dalam aljabar-$*\,$}%
 \index{<star@homomorfisme-$*\,$!beridentitas}%
\df{beridentitas} jika $\phi(\vc 1_A) = \vc 1_B$.  Dalam kategori aljabar-$*\,$ dan homomorfisme-$*\,$, isomorfisme (yang untuk penegasan disebut
 \index{<star@isomorfisme-$*\,$}%
 \index{star@isomorfisme-$*\,$}%
 \index{isomorfisme!$*\,$-}
\df{isomorfisme-$*\,$}) adalah homomorfisme-$*\,$ bijektif.
\end{defn}

\begin{prop}\label{0010952} Misalkan $\phi\colon A \sto B$ suatu homomorfisme-$*\,$ di antara aljabar-$*\,$. Maka
kernel $\phi$ adalah suatu ideal-$*\,$ dalam $A$, dan citra $\phi$ adalah suatu
subaljabar-$*\,$ dari~$B$.
\end{prop}
















\section{Operator Swaadjoin}
\begin{defn} Suatu elemen $a$ dari aljabar-$*\,$ $A$ bersifat
 \index{swaadjoin!elemen}%
\df{swaadjoin} (atau
 \index{Hermitian!elemen}%
\df{Hermitian}) jika $a^* = a$.  Elemen ini bersifat
 \index{normal!elemen}%
\df{normal} jika $a^*a = aa^*$. Jika aljabarnya beridentitas, elemen ini bersifat
 \index{uniter!elemen}%
\df{uniter} jika $a^*a = aa^* = \vc 1$. Himpunan semua elemen swaadjoin dari $A$
dinotasikan dengan
 \index{H@$\fml H(A)$ (elemen swaadjoin dari suatu aljabar-$*\,$)}%
$\fml H(A)$, himpunan semua elemen normal
 \index{N@$\fml N(A)$ (elemen normal dari suatu aljabar-$*\,$)}%
dengan~$\fml N(A)$, dan, ketika aljabarnya beridentitas, himpunan semua elemen uniter
 \index{U@$\fml U(A)$ (elemen uniter dari suatu aljabar-$*\,$)}%
dengan~$\fml U(A)$.
\end{defn}

\begin{prop}\label{001093} Misalkan $a$ elemen dari suatu aljabar-$*\,$ kompleks.  Maka terdapat
elemen swaadjoin unik $u$ dan $v$ sedemikian sehingga $a = u + iv$.
\end{prop}

\begin{proof}[\emph{Petunjuk pembuktian}] Pikirkan kasus khusus penulisan suatu bilangan kompleks
dalam bentuk bagian real dan imajinernya.  \ns
\end{proof}

\begin{defn} Ambil $S$ suatu himpunan bagian dari aljabar-$*\,$~$A$.  Maka
$S^*~=~\{s^*\colon~s~\in~S\}$. Himpunan $S$ disebut
 \index{swaadjoin!himpunan bagian dari suatu aljabar-$*\,$}%
 \index{<unopurstar@$S^*$ (himpunan adjoin dari elemen-elemen~$S$)}%
\df{swaadjoin} jika $S^* = S$.

Suatu subaljabar swaadjoin tak kosong dari $A$ disebut
 \index{<star@subaljabar-$*\,$ (subaljabar swaadjoin)}%
 \index{star@subaljabar-$*\,$ (subaljabar swaadjoin)}%
 \index{subaljabar!$*\,$-}%
 \index{swaadjoin!subaljabar}%
\df{subaljabar-$*\,$}. Istilah sumber yang bersinonim juga diterjemahkan sebagai \df{subaljabar-$*\,$}.
\end{defn}

\begin{cau} Definisi sebelumnya \emph{tidak} menyatakan bahwa elemen-elemen dari suatu
himpunan bagian swaadjoin dari suatu aljabar-$*\,$ dengan sendirinya bersifat swaadjoin.
\end{cau}

 \begin{prop} Suatu operator $T$ pada ruang Hilbert kompleks $H$ bersifat swaadjoin jika dan hanya jika
 bentuk kuadratik terkaitnya $Q_T$ bernilai real.
 \end{prop}

\begin{proof}[\emph{Petunjuk pembuktian}] Gunakan trik yang sama seperti dalam~\ref{C0635105}.  Dalam hipotesis bahwa
$Q_T(z)$ selalu real, gantilah $z$ mula-mula dengan $y + x$, lalu dengan $y + ix$.   \ns
\end{proof}

\begin{defn} Misalkan $T$ suatu operator pada ruang Hilbert tak nol~$H$.  Adapun
 \index{numerik!jangkauan}%
 \index{jangkauan!numerik}%
 \index{W@$W(T)$ (jangkauan numerik)}%
\df{jangkauan numerik} dari $T$, yang dinotasikan dengan $W(T)$, adalah jangkauan yang diperoleh pada sfer satuan dari $H$
dengan membatasi bentuk kuadratik yang terkait dengan~$T$. Yaitu,
   \[ W(T) := \{Q_T(x) \colon x \in H \text{ dan } \norm x = 1\}. \]
Adapun
 \index{numerik!radius}%
 \index{radius!numerik}%
 \index{w@$w(T)$ (radius numerik)}%
\df{radius numerik} dari $T$, yang dinotasikan dengan $w(T)$, adalah $\sup\{\abs \lambda\colon \lambda \in W(T)\}$.
\end{defn}

\begin{prop} Jika $T$ suatu operator swaadjoin pada ruang Hilbert tak nol, maka $\norm T = w(T)$.
\end{prop}

\begin{proof}[\emph{Petunjuk pembuktian}] Bagian yang sulit adalah menunjukkan bahwa $\norm T \le w(T)$.  Untuk itu, tunjukkan terlebih dahulu bahwa
   \begin{equation}
             Q_T(x + y) - Q_T(x - y) =   4 \Re\langle Tx,y \rangle
   \end{equation}
untuk semua vektor $x$ dan $y$ dalam ruang Hilbert, di mana $Q_T$ merupakan bentuk kuadratik yang terkait dengan $T$ dan $\Re z$ adalah
bagian real dari bilangan kompleks~$z$.  Gunakan ini untuk memverifikasi bahwa jika $\norm x = \norm y = 1$, maka
   \begin{equation}\label{num_ran_saop_eqn2}
           \Re\langle Tx,y \rangle \le w(T).
   \end{equation}
Tuliskan bilangan kompleks $\langle Tx,y \rangle$ dalam bentuk polar $re^{i\theta}$.  Jelas bahwa jika kita mengganti
$x$ dengan $e^{-i\theta}x$ dalam pertidaksamaan~\eqref{num_ran_saop_eqn2}, pertidaksamaan yang dihasilkan tetap berlaku kapan pun $\norm x = \norm y = 1$.

Terakhir, perhatikan bahwa jika $\norm x = 1$ dan $y$ adalah sebarang vektor dalam $H$ sedemikian sehingga $Tx = \norm{Tx}\,y$, maka
   \begin{equation}
             \norm{Tx} = \langle Tx,y \rangle\,.
   \end{equation}
Dengan demikian, hasil yang diinginkan segera diperoleh.         \ns
\end{proof}

Dalam proposisi~\ref{C0635105}, kita melihat bahwa suatu operator $T$ pada ruang Hilbert \emph{kompleks} merupakan operator nol jika dan hanya jika
bentuk kuadratik terkaitnya bernilai nol.  Dari proposisi sebelumnya mudah diperoleh akibat bahwa ekuivalensi ini juga berlaku
untuk operator swaadjoin pada ruang Hilbert real.

\begin{cor}\label{real_HSp_zero_Qform} Misalkan $T$ suatu operator swaadjoin pada ruang Hilbert.  Maka $T = \vc 0$ jika dan
hanya jika bentuk kuadratik terkaitnya $Q_T$ bernilai nol.
\end{cor}

\begin{defn} Suatu operator $T$ pada ruang Hilbert $H$ bersifat
 \index{positif!operator}%
 \index{operator!positif}%
\df{positif} jika operator tersebut swaadjoin dan bentuk kuadratik terkaitnya positif ($Q_T(x) \ge 0$ untuk setiap $x \in H$).
\end{defn}

\begin{exer} Misalkan $H$ ruang Hilbert kompleks dan $T$ operator positif pada~$H$.  Definisikan
   \[ \langle x,y \rangle_1 := \langle Tx,y \rangle \qquad \text{ dan } \qquad \norm{x}_1 \equiv \sqrt{\langle x,x \rangle_1} \]
untuk semua $x \in H$.
 \begin{enumerate}
  \item[(a)] Buktikan bahwa $\norm{\hphantom O}_1$ merupakan norma jika dan hanya jika $\ker T = \{0\}$.
  \item[(b)] Misalkan $\ker T = \{0\}$. Tunjukkan bahwa norma $\norm{\hphantom O}$ dan
$\norm{\hphantom O}_1$ menginduksi topologi yang sama pada~$H$ jika dan hanya jika $T$ invertibel dalam
$\ofml B(H)$. \emph{Petunjuk.}  Ketika terdapat dua topologi pada~$H$, sering kali berguna untuk
mempertimbangkan operator identitas pada~$H$.
  \item[(c)] Misalkan $T$ invertibel dalam $\ofml B(H)$.  Untuk $S \in \ofml B(H)$, carilah suatu rumus bagi
adjoin dari $S$ terhadap hasil kali dalam $\langle \hphantom x , \hphantom y \rangle_1$.
 \end{enumerate}
\end{exer}


















\section{Proyeksi}
\begin{conv} Saat ini perhatian kita terutama difokuskan pada
 \index{konvensi!proyeksi di ruang Hilbert bersifat ortogonal}%
operator pada ruang Hilbert.  Dalam konteks ini, istilah \emph{proyeksi} selalu dimaknai
 \index{proyeksi}%
sebagai \emph{proyeksi ortogonal} (lihat definisi~\ref{000164}). Jadi, suatu operator pada ruang Hilbert
disebut proyeksi jika $P^2 = P$ dan $P^* = P$.
\end{conv}

Kita memperumum definisi ``proyeksi (ortogonal)'' dari ruang Hilbert ke aljabar-$*\,$.

\begin{defn} Suatu
 \index{proyeksi!dalam suatu aljabar-$*\,$}%
\df{proyeksi} dalam suatu aljabar-$*\,$ $A$ adalah suatu elemen $p$ dari aljabar tersebut yang bersifat
idempoten ($p^2 = p$) dan swaadjoin ($p^* = p$).  Himpunan semua proyeksi dalam $A$
dinotasikan
 \index{P@$\fml P(A)$!proyeksi dalam suatu aljabar}%
dengan~$\fml P(A)$.  Dalam kasus $\ofml B(H)$, yaitu operator terbatas pada suatu ruang Hilbert,
kita menulis $\fml P(H)$,
 \index{p@$\fml P(H)$ atau $\fml P(\ofml B(H))$!proyeksi ortogonal pada ruang Hilbert}%
atau, jika $H$ telah dipahami, cukup $\fml P$, untuk notasi yang lebih informatif $\fml P(\ofml B(H))$.
\end{defn}

\begin{prop} Setiap operator pada ruang Hilbert yang merupakan isometri pada komplemen ortogonal dari kernelnya memiliki jangkauan tertutup.
\end{prop}

\begin{proof}[\emph{Petunjuk pembuktian}] Perhatikan terlebih dahulu bahwa jika $M = (\ker T)^\perp$, dengan $T$ adalah operator yang dimaksud, maka
$\ran T = T^{\sto}(M)$.    \ns
\end{proof}

\begin{prop} Misalkan $P$ suatu proyeksi pada ruang Hilbert $H$. Maka
 \begin{enumerate}[label=\rm{(\alph*)}]
  \item $Px = x$ jika dan hanya jika $x \in \ran P$\,;
  \item $\ker P = (\ran P)^\perp$; dan
  \item $H = \ker P \oplus \ran P$.
 \end{enumerate}
\end{prop}

Pada bagian (c) dari proposisi sebelumnya, lambang $\oplus$ tentu saja menyatakan
jumlah langsung \emph{ortogonal} (lihat paragraf setelah~\ref{000164}).

\begin{prop} Misalkan $M$ suatu subruang dari ruang Hilbert $H$.  Jika $P \in \ofml B(H)$, jika $Px = x$
untuk setiap $x \in M$, dan jika $Px = \vc 0$ untuk setiap $x \in M^\perp$, maka $P$ adalah
proyeksi dari $H$ pada~$M$.
\end{prop}

Pada bagian selebihnya dari seksi ini, penting untuk selalu mengingat bahwa ruang $\ofml B(H)$ dari operator pada ruang
Hilbert merupakan \emph{contoh} aljabar-$*\,$, dan bahwa proyeksi (ortogonal) pada suatu ruang Hilbert merupakan
\emph{contoh} proyeksi dalam aljabar-$*\,$.  Jadi, di satu sisi, semua hal yang dinyatakan untuk
proyeksi semacam itu dalam aljabar-$*\,$ berlaku bagi proyeksi ruang Hilbert. Karena tidak ada nama yang lebih baik, kita akan mengatakan bahwa
hasil-hasil ini mendokumentasikan struktur
 \index{pandangan abstrak tentang proyeksi}%
\df{abstrak} dari proyeksi.  Di sisi lain, proyeksi ruang Hilbert memiliki struktur tambahan yang
tidak bermakna dalam aljabar-$*\,$ umum: proyeksi tersebut memiliki \emph{domain} dan \emph{jangkauan}, serta \emph{bekerja pada vektor}.
Ketika menyelidiki hal-hal semacam itu, kita akan mengatakan bahwa kita membahas struktur
 \index{pandangan spasial tentang proyeksi}%
\df{spasial} (atau \df{konkret}) dari proyeksi.  Perhatikan bagaimana kedua pandangan ini muncul berselang-seling dalam bagian selebihnya dari seksi ini.

Hasil pertama kita memberikan syarat perlu dan cukup agar jumlah proyeksi merupakan suatu proyeksi.

\begin{prop}\label{prop_sum_projs} Misalkan $p$ dan $q$ proyeksi dalam suatu aljabar-$*\,$. Maka pernyataan berikut ekuivalen:
 \begin{enumerate}[label=\rm{(\alph*)}]
  \item $pq = \vc 0$;
  \item $qp = \vc 0$;
  \item $qp = -pq$;
 \index{proyeksi!jumlah dari}%
  \item $p + q$ merupakan suatu proyeksi.
 \end{enumerate}
\end{prop}

\begin{defn}\label{def_perp_projs} Misalkan $p$ dan $q$ proyeksi dalam suatu aljabar-$*\,$.  Jika salah satu syarat
dalam hasil sebelumnya berlaku, maka kita mengatakan bahwa $p$ dan $q$ bersifat
 \index{proyeksi!ortogonal}%
 \index{proyeksi!keortogonalan dari dua proyeksi}%
 \index{ortogonal!proyeksi!dalam suatu aljabar-$*\,$}%
\df{ortogonal} dan menulis
 \index{<binrel@$p \perp q$ (keortogonalan proyeksi)}%
$p \perp q$.  (Jadi, untuk operator pada suatu ruang Hilbert, kita dapat dengan tepat, meskipun tidak nyaman, menyebut
\emph{proyeksi ortogonal yang ortogonal}!)
\end{defn}

Berikutnya adalah hasil spasial yang bersesuaian untuk jumlah dua proyeksi.

\begin{prop} Misalkan $P$ dan $Q$ proyeksi pada ruang Hilbert~$H$. Maka $P \perp Q$
 \index{ortogonal!proyeksi ruang Hilbert}%
jika dan hanya jika $\ran P \perp \ran Q$. Dalam kasus ini, $P + Q$ merupakan suatu proyeksi
yang kernelnya adalah $\ker P \cap \ker Q$ dan jangkauannya adalah $\ran P + \ran Q$.
\end{prop}

\begin{exam} Pada suatu ruang Hilbert, proyeksi (ortogonal) tidak harus berkomutasi. Sebagai contoh,
misalkan $P$ proyeksi dari ruang Hilbert (real) $\R^2$ pada garis $y = x$ dan $Q$
adalah proyeksi dari $\R^2$ pada sumbu-$x$. Maka $PQ \ne QP$.
\end{exam}

Hasil kali dua proyeksi merupakan suatu proyeksi jika dan hanya jika keduanya berkomutasi.

\begin{prop} Misalkan $p$ dan $q$ proyeksi dalam suatu aljabar-$*\,$. Maka $pq$ merupakan suatu
 \index{proyeksi!hasil kali dari}%
proyeksi jika dan hanya jika $pq = qp$.
\end{prop}

\begin{prop} Misalkan $P$ dan $Q$ proyeksi pada ruang Hilbert~$H$. Jika $PQ = QP$, maka
$PQ$ merupakan suatu proyeksi yang kernelnya adalah $\ker P + \ker Q$ dan jangkauannya adalah $\ran P \cap
\ran Q$.
\end{prop}

\begin{prop} Misalkan $p$ dan $q$ proyeksi dalam suatu aljabar-$*\,$. Maka pernyataan berikut ekuivalen:
 \begin{enumerate}[label=\rm{(\alph*)}]
  \item $pq = p$;
  \item $qp = p$;
 \index{proyeksi!selisih dari}%
  \item $q - p$ merupakan suatu proyeksi.
 \end{enumerate}
\end{prop}

\begin{defn}\label{0022251} Misalkan $p$ dan $q$ proyeksi dalam suatu aljabar-$*\,$.  Jika salah satu
dari syarat dalam hasil sebelumnya berlaku, maka kita
 \index{<binrelle@$p \preceq q$ (urutan proyeksi dalam suatu aljabar-$*\,$)}%
 \index{proyeksi!urutan dari}%
 \index{urutan!parsial!dari proyeksi}%
 \index{parsial!urutan!dari proyeksi}%
menulis $p \preceq q$.
\end{defn}

\begin{prop} Misalkan $P$ dan $Q$ proyeksi pada ruang Hilbert~$H$. Maka pernyataan berikut ekuivalen:
 \begin{enumerate}[label=\rm{(\alph*)}]
  \item $P \preceq Q$;
  \item $\norm{Px} \le \norm{Qx}$ untuk semua $x \in H$; dan
  \item $\ran P \subseteq \ran Q$.
 \end{enumerate}
Dalam kasus ini, $Q - P$ merupakan suatu proyeksi yang kernelnya adalah $\ran P + \ker Q$ dan jangkauannya adalah $\ran Q \ominus \ran P$.
\end{prop}

\noindent\emph{Notasi:} Misalkan $H$, $M$, dan $N$ subruang dari suatu ruang Hilbert. Pernyataan
 \index{<binopdirectdiff@$\ominus$ (selisih langsung)}%
$H = M \oplus N$ dapat ditulis ulang sebagai $M = H \ominus N$ (atau $N = H \ominus M$).

\begin{prop} Operasi $\preceq$ yang didefinisikan dalam~\ref{0022251} untuk proyeksi pada suatu
aljabar-$*\,$ beridentitas $A$ merupakan urutan parsial pada $\fml P(A)$. Jika $p$ adalah suatu proyeksi dalam $A$,
maka $\vc 0 \preceq p \preceq \vc 1$.
\end{prop}

\begin{prop} Misalkan $p$ dan $q$ proyeksi dalam suatu aljabar-$*\,$~$A$. Jika $pq = qp$, maka
infimum dari $p$ dan $q$, yang kita notasikan
 \index{<binrel@$p \curlywedge q$ (infimum proyeksi)}%
 \index{infimum!dari proyeksi dalam suatu aljabar-$*\,$}%
 \index{proyeksi!infimum dari}%
dengan $p \curlywedge q$, ada terhadap urutan parsial $\preceq$ dan $p
\curlywedge q = pq$.  Infimum $p \curlywedge q$ dapat ada bahkan ketika $p$ dan $q$ tidak
berkomutasi.  Syarat perlu dan cukup agar $p \perp q$ berlaku adalah bahwa $p
\curlywedge q = \vc 0$ dan $pq = qp$ keduanya berlaku.
\end{prop}

\begin{prop} Misalkan $p$ dan $q$ proyeksi dalam suatu aljabar-$*\,$~$A$. Jika $p \perp q$, maka
supremum dari $p$ dan $q$, yang kita notasikan
 \index{<binrel@$p \curlyvee q$ (supremum proyeksi)}%
 \index{supremum!dari proyeksi dalam suatu aljabar-$*\,$}%
 \index{proyeksi!supremum dari}%
dengan $p \curlyvee q$, ada terhadap urutan parsial $\preceq$ dan $p
\curlyvee q = p + q$.  Supremum $p \curlyvee q$ dapat ada bahkan ketika $p$ dan $q$ tidak
ortogonal.
\end{prop}

\begin{prop} Misalkan $A$ suatu aljabar-$C^*$, $a \in A$ suatu elemen positif, dan $0 < \epsilon \le \frac12$. Jika
$\norm{a^2 - a} < \epsilon/2$, maka terdapat suatu proyeksi $p$ dalam $A$ sedemikian sehingga
$\norm{p - a} < \epsilon$.
\end{prop}
 
 
 
 
 



 
 

\section{Operator Normal}

\begin{prop}\label{prop_nasc_normal} Suatu operator $T$ pada ruang Hilbert $H$ bersifat normal jika dan hanya jika
$\norm{Tx} = \norm{T^*x}$ untuk semua $x \in H$.
\end{prop}

\begin{proof}[\emph{Petunjuk pembuktian}] Gunakan Korolar~\ref{real_HSp_zero_Qform}.  \ns
\end{proof}

\begin{cor} Jika $T$ merupakan operator normal pada suatu ruang Hilbert, maka $\ker T = \ker T^*$.
\end{cor}

\begin{prop} Jika $T$ merupakan operator normal pada suatu ruang Hilbert, maka $\norm{T^2} = {\norm T}^2$.
\end{prop}

\begin{proof}[\emph{Petunjuk pembuktian}] Substitusikan $Tx$ sebagai pengganti $x$ dalam Proposisi~\ref{prop_nasc_normal}. Gunakan
Proposisi~\ref{C063544}.  \ns
\end{proof}

\begin{prop} Jika $T$ merupakan operator normal pada ruang Hilbert $H$, maka
   \[ H = \clo{\ran T} \oplus \ker T\,.  \]
\end{prop}

\begin{cor} Suatu operator normal pada ruang Hilbert memiliki jangkauan padat jika dan hanya jika operator itu injektif.
\end{cor}

\begin{exam} Operator \emph{geser unilateral} (\ref{C063526}) merupakan contoh operator pada ruang Hilbert yang injektif
tetapi tidak memiliki jangkauan padat. Adjoinnya memiliki jangkauan padat (bahkan, bersifat surjektif), tetapi tidak injektif.
\end{exam}

\begin{cor} Suatu operator normal pada ruang Hilbert bersifat invertibel jika dan hanya jika operator tersebut terbatas jauh dari nol.
\end{cor}
















\section{Operator Berperingkat Hingga}

\begin{defn}
 \index{peringkat}%
\df{peringkat} suatu pemetaan linear adalah dimensi (ruang vektor) dari jangkauannya. Jadi, suatu operator
 \index{hingga!peringkat}%
 \index{operator!berperingkat hingga}%
\df{berperingkat hingga} adalah operator yang memiliki jangkauan berdimensi hingga. Kita menyatakan dengan $\ofml{FR}(V)$
 \index{fr@$\ofml{FR}(V)$ (operator berperingkat hingga pada~$V$)}%
keluarga operator berperingkat hingga pada suatu ruang vektor $V$.
\end{defn}

\begin{exam}\label{0022431fa} Untuk vektor $x$ dan $y$ dalam ruang Hilbert $H$, definisikan
   \[ x \otimes y\colon H \sto H \colon z \mapsto \langle z,y \rangle x\,. \]
Jika $x$ dan $y$ merupakan vektor tak nol dalam $H$, maka $x \otimes y$ merupakan operator berperingkat satu pada~$H$.
\end{exam}

Dua proposisi berikutnya mengidentifikasi adjoin dan komposisi operator-operator berperingkat satu ini.

\begin{prop} Jika $u$ dan $v$ merupakan vektor dalam suatu ruang Hilbert, maka
   \[ (u \otimes v)^* = v \otimes u\,. \]
\end{prop}

\begin{prop} Jika $u$, $v$, $x$, dan $y$ merupakan vektor dalam suatu ruang Hilbert, maka
    \[ (u \otimes v)(x \otimes y) = \langle x,v \rangle (u \otimes y)\,. \]
\end{prop}

\begin{prop} Jika $x$ merupakan vektor dalam ruang Hilbert $H$, maka $x \otimes x$ merupakan proyeksi berperingkat satu jika dan
hanya jika $x$ merupakan vektor satuan.
\end{prop}

\begin{prop} Andaikan bahwa $\{e^1, e^2, e^3, \dots, e^n\}$ merupakan himpunan ortonormal dalam ruang Hilbert~$H$. Misalkan
$M = \spn\{e^1, \dots, e^n\}$. Maka $P = \sum_{k=1}^n e^k \otimes e^k$ merupakan proyeksi ortogonal dari $H$ ke~$M$.
\end{prop}

\begin{prop}\label{prop_op_act_otimes1} Jika $u$ dan $v$ merupakan vektor dalam suatu ruang Hilbert dan $T$ merupakan operator pada ruang Hilbert yang sama, yaitu~$H$, maka
   \[ T\,(u \otimes v) = (Tu) \otimes v\,. \]
\end{prop}

\begin{prop}\label{prop_op_act_otimes2} Jika $u$ dan $v$ merupakan vektor dalam suatu ruang Hilbert dan $T$ merupakan operator pada ruang Hilbert yang sama, yaitu~$H$, maka
   \[ (u \otimes v)\,T = u \otimes T^*v\,. \]
\end{prop}

Kita akan segera melihat bahwa keluarga semua operator berperingkat hingga pada suatu ruang Hilbert tak nol merupakan ideal-$*\,$ minimal dalam
aljabar-$*\,$ $\ofml B(H)$. Sebagai persiapan, kita meninjau beberapa fakta dasar tentang ideal dalam aljabar.

\begin{defn}\label{000551} Suatu
 \index{ideal!kiri}%
 \index{kiri!ideal}%
\df{ideal kiri} dalam suatu aljabar $A$ adalah subruang vektor $J$ dari $A$ sedemikian sehingga $AJ \subseteq J$. (Untuk
 \index{ideal!kanan}%
 \index{kanan!ideal}%
\df{ideal kanan}, tentu saja, kita mensyaratkan $JA \subseteq J$.) Kita mengatakan bahwa $J$ merupakan suatu
 \index{ideal!dalam suatu aljabar}%
\df{ideal} jika ideal tersebut merupakan ideal dua sisi, yaitu ideal kiri sekaligus ideal kanan. Suatu ideal disebut
 \index{sejati!ideal}%
 \index{ideal!sejati}%
\df{sejati} jika ideal itu merupakan subhimpunan sejati dari~$A$.

Ideal $\{0\}$ dan $A$ sering disebut sebagai
 \index{trivial!ideal}%
 \index{ideal!trivial}%
\df{ideal trivial} dari~$A$. Aljabar $A$ disebut
 \index{sederhana!aljabar}%
 \index{aljabar!sederhana}%
\df{sederhana} jika tidak memiliki ideal tak trivial.

Suatu
 \index{ideal maksimal}%
 \index{ideal!maksimal}%
ideal \df{maksimal} adalah ideal sejati yang tidak termuat secara sejati dalam ideal sejati lain.
Kita menyatakan keluarga semua ideal maksimal dalam suatu aljabar $A$
 \index{maximal@$\mx A$!himpunan ideal maksimal dalam suatu aljabar}%
dengan~$\mx A$. Suatu
 \index{ideal minimal}%
 \index{ideal!minimal}%
ideal \df{minimal} adalah ideal tak nol yang tidak memuat secara sejati ideal tak nol lain.
\end{defn}

 \index{konvensi!ideal bersifat dua sisi}%
\begin{conv} Setiap kali kita merujuk pada suatu \emph{ideal} dalam aljabar, ideal tersebut dipahami sebagai
ideal dua sisi (kecuali dinyatakan sebaliknya).
\end{conv}

\begin{prop} Tidak satu pun elemen invertibel dalam aljabar beridentitas dapat termasuk dalam suatu ideal sejati.
\end{prop}

 \begin{prop}\label{000552} Setiap ideal sejati dalam aljabar beridentitas $A$ termuat dalam suatu
 ideal maksimal. Jadi, khususnya, $\mx A$ tak kosong apabila $A$ merupakan aljabar beridentitas.
\end{prop}

\begin{proof}[\emph{Petunjuk pembuktian}] Lema Zorn. \ns  \end{proof}

\begin{prop}\label{000553} Misalkan $a$ suatu elemen dari aljabar komutatif~$A$. Jika $A$
beridentitas dan $a$ tidak invertibel, maka $aA$ merupakan ideal sejati dalam~$A$.
\end{prop}

\begin{defn}\label{000554} Ideal $aA$ dalam proposisi sebelumnya adalah
 \index{ideal utama}%
 \index{ideal!utama}%
\df{ideal utama} yang dibangkitkan oleh~$a$.
\end{defn}

\begin{defn}\label{000555} Misalkan $J$ suatu ideal dalam aljabar~$A$. Definisikan suatu
relasi ekuivalensi $\sim$ pada $A$ dengan
  \[ a \sim b \qquad \text{jika dan hanya jika } \qquad b-a \in J. \]
Untuk setiap $a \in A$, misalkan $[a]$ adalah kelas ekuivalensi yang memuat~$a$. Misalkan $A/J$ adalah
himpunan semua kelas ekuivalensi elemen-elemen~$A$. Untuk $[a]$ dan $[b]$ dalam $A/J$, definisikan
  \[ [a] + [b] := [a+b] \qquad \text{ dan } \qquad [a][b] := [ab] \]
dan untuk $\alpha \in \C$ dan $[a] \in A/J$, definisikan
  \[\alpha[a] := [\alpha a] \,. \]
Terhadap operasi-operasi ini, $A/J$ menjadi suatu aljabar. Aljabar tersebut adalah
 \index{hasil bagi!aljabar}%
 \index{aljabar!hasil bagi}%
\df{aljabar hasil bagi} dari $A$ oleh~$J$. Notasi $A/J$ biasanya dibaca ``$A$ modulo~$J$''.
Homomorfisme aljabar surjektif
  \[ \pi\colon A \sto A/J\colon a \mapsto [a]\]
disebut
 \index{hasil bagi!pemetaan!untuk aljabar}%
\df{pemetaan hasil bagi}.
\end{defn}

\begin{prop} Definisi sebelumnya bermakna dan pernyataan-pernyataan di dalamnya benar.
Lebih lanjut, jika ideal $J$ bersifat sejati, maka aljabar hasil bagi $A/J$ beridentitas jika $A$ beridentitas.
\end{prop}

\begin{proof}[\emph{Petunjuk pembuktian}] Anda perlu menunjukkan bahwa:
 \begin{enumerate}
  \item[(a)] $\sim$ merupakan relasi ekuivalensi.
  \item[(b)] Penjumlahan dan perkalian kelas-kelas ekuivalensi terdefinisi dengan baik.
  \item[(c)] Perkalian suatu kelas ekuivalensi dengan skalar terdefinisi dengan baik.
  \item[(d)] $A/J$ merupakan suatu aljabar.
  \item[(e)] ``Pemetaan hasil bagi'' $\pi$ memang merupakan homomorfisme aljabar surjektif.
 \end{enumerate}
(Pada tahap mana diperlukan bahwa kita mengambil hasil bagi oleh suatu ideal, bukan sekadar subaljabar?) \ns
\end{proof}

\begin{defn} Dalam suatu aljabar $A$ dengan involusi,
 \index{<star@ideal-$*\,$}%
 \index{star@ideal-$*\,$}%
 \index{ideal!$*\,$-}%
suatu \df{ideal-$*\,$} adalah ideal swaadjoin dalam~$A$.
\end{defn}

\begin{prop}\label{0010953} Andaikan $J$ suatu ideal-$*\,$ dalam aljabar-$*\,$ $A$. Dalam aljabar
hasil bagi $A/J$, sebagaimana dikembangkan dalam~\ref{000555}, definisikan $[a]^* := [a^*]$ untuk setiap $[a] \in A/J$.
 \index{hasil bagi!aljabar-$*\,$}%
 \index{<star@aljabar-$*\,$!hasil bagi}%
Ini merupakan operasi uner yang terdefinisi dengan baik pada aljabar hasil bagi $A/J$ yang terbentuk; aljabar tersebut menjadi aljabar-$*\,$, dan
pemetaan hasil bagi $a \mapsto [a]$ merupakan homomorfisme-$*\,$. Aljabar-$*\,$ $A/J$ adalah, tentu
saja,
 \index{hasil bagi!aljabar-$*\,$}%
 \index{<star@aljabar-$*\,$!hasil bagi}%
 \index{star@aljabar-$*\,$!hasil bagi}%
\df{hasil bagi} dari $A$ oleh~$J$.
\end{prop}

\begin{prop} Misalkan $H$ suatu ruang Hilbert. Maka keluarga $\ofml{FR}(H)$ dari semua operator berperingkat hingga pada $H$
 \index{star@ideal-$*\,$!$\ofml{FR}(H)$ sebagai suatu}%
 \index{FR@$\ofml{FR}(H)$!sebagai ideal-$*\,$}%
merupakan ideal-$*\,$ dalam aljabar-$*\,$~$\ofml B(H)$.
\end{prop}

\begin{prop}\label{prop_frop_sum_rnk1} Misalkan $T$ suatu operator berperingkat $n$ pada ruang Hilbert~$H$. Maka terdapat
vektor $u^1$, \dots, $u^n$, $v^1$, \dots, $v^n$ dalam $H$ sedemikian sehingga
   \[ T = \sum_{k=1}^n u^k \otimes v^k\,. \]
\end{prop}

\begin{proof}[\emph{Petunjuk pembuktian}] Misalkan $\{e^1, \dots, e^n\}$ suatu basis ortonormal bagi jangkauan~$T$. Untuk
sembarang vektor $x \in H$, tuliskan \emph{ekspansi Fourier} dari $Tx$ (lihat Proposisi~\ref{C063144}(d)\,).  \ns
\end{proof}

\begin{cor} Setiap operator berperingkat satu pada suatu ruang Hilbert berbentuk $u \otimes v$ untuk suatu pasangan vektor tak nol
$u$ dan $v$ dalam~$H$.
\end{cor}

\begin{cor} Setiap operator berperingkat hingga pada suatu ruang Hilbert merupakan jumlah dari (sejumlah berhingga) operator berperingkat satu.
\end{cor}

\begin{prop} Jika $H \ne \{0\}$ merupakan ruang Hilbert, maka keluarga $\ofml{FR}(H)$ dari operator berperingkat hingga pada $H$
merupakan ideal minimal dalam~$\ofml B(H)$.
\end{prop}

\begin{proof}[\emph{Petunjuk pembuktian}] Misalkan $\ofml J$ suatu ideal tak nol dalam $\ofml B(H)$. Kita perlu menunjukkan bahwa $\ofml J$
memuat~$\ofml{FR}(H)$. Berdasarkan Proposisi~\ref{prop_frop_sum_rnk1}, cukup dibuktikan bahwa $\ofml J$ memuat setiap
operator berbentuk $u \otimes v$. Untuk itu, misalkan $T$ sebarang anggota tak nol dari $\ofml J$, misalkan $y$ suatu vektor satuan dalam
jangkauan $T$, dan misalkan $x$ sebarang vektor dalam $H$ sedemikian sehingga $y = Tx$. Tinjau operator $(u \otimes y)T(x \otimes v)$. \ns
\end{proof}

Meskipun himpunan operator berperingkat hingga merupakan ideal dalam aljabar Banach yang terdiri atas operator pada suatu ruang Hilbert berdimensi tak hingga, himpunan itu bukan ideal
tertutup dan, sebagai akibatnya, ternyata bukan himpunan yang tepat untuk membentuk hasil bagi dengan harapan memperoleh suatu objek hasil bagi dalam
kategori aljabar Banach. Tutupan ideal operator berperingkat hingga pada suatu ruang Hilbert adalah ideal \emph{operator
kompak}. Dalam arti tertentu, langkah berikutnya yang wajar ialah mengkaji sifat-sifat kelas operator ini; di dalamnya kita menjumpai
beberapa operator terpenting dalam penerapan, misalnya operator integral. Namun, pemaparan yang lebih lancar dan lebih gamblang
dimungkinkan jika kita memiliki sedikit lebih banyak perangkat teknis. Oleh karena itu, dalam bab berikutnya kita akan berfokus pada
pengembangan beberapa alat baku yang digunakan dalam analisis fungsional.
























\endinput
